\documentclass[11pt,a4paper]{article}
\usepackage[T1]{fontenc}
\usepackage[utf8]{inputenc}
\usepackage{mathpazo}
\usepackage[scaled=0.92]{helvet}

\usepackage{microtype}
\usepackage{graphicx}
\usepackage[font=small,labelfont=bf]{caption}
\renewcommand{\thefigure}{12.\arabic{figure}}
\usepackage[margin=2.6cm]{geometry}
\usepackage[round]{natbib}
\usepackage{titlesec}
\usepackage{enumitem}
\usepackage{xcolor}
\usepackage[most]{tcolorbox}
\usepackage[hidelinks]{hyperref}
\usepackage{url}
\urlstyle{same}

\setlist[itemize]{itemsep=2pt,topsep=4pt}
\setlength{\parskip}{6pt}
\setlength{\parindent}{0pt}
\renewcommand{\thesection}{12.\arabic{section}}

% Footnote citation in the style of the first edition
\newcommand{\fn}[1]{\footnote{\citealp{#1}.}}

% Claim-status marks (underlines, margin bars and the key)
\usepackage{gpt-status}
\usepackage{longtable}
\usepackage{array}
\usepackage{booktabs}
\usepackage{float}
\usepackage{placeins}
% Explanatory registers
\newcommand{\reg}[1]{\textsf{\small\textbf{[#1]}}}
\newcommand{\RN}{\reg{N}}
\newcommand{\RGR}{\reg{GR}}
\newcommand{\RFR}{\reg{FR}}
\newcommand{\RC}{\reg{C}}
\newcommand{\REN}{\reg{EN}}
\newcommand{\RQ}{\reg{Q}}
\newcommand{\RDS}{\reg{DS}}
\newcommand{\expl}[1]{\par\noindent\hspace*{1.2em}\begin{minipage}[t]{\dimexpr\linewidth-1.2em}\small\textit{Explained as:} #1\end{minipage}\par\smallskip}
\newcommand{\why}[1]{\par\textit{#1}\par}
\usepackage{tocloft}
\renewcommand{\cftsecfont}{\bfseries}\renewcommand{\cftsecpagefont}{\bfseries}\setlength{\cftsecnumwidth}{3.2em}
\begin{document}
\renewcommand{\thetable}{12.\arabic{table}}
\begin{titlepage}
\centering
\vspace*{1.5cm}
{\Huge\bfseries Graviton Pressure Theory\par}
\vspace{0.4cm}
{\Large The Unified Framework\par}
{\large Individual Submission\par}
\vspace{0.3cm}
{\normalsize This document is part of a multi-part scientific framework\par}
\vspace{1.6cm}
{\large Part 12 of 30\par}
\vspace{0.6cm}
{\huge\bfseries The Collapse of Materialist Epistemology\par}
\vspace{0.5cm}
{\Large Exposing the Logical and Scientific Limits of an Unravelling Paradigm\par}
\vspace{1.4cm}
\begin{minipage}{0.82\textwidth}\centering\small
This submission is part of the broader Graviton Pressure Theory (GPT) project, a redefinition of gravitational interaction rooted in causal field dynamics and physical force transmission. Each document is designed to stand on its own; its full context emerges as part of the larger framework. For the complete series, see \url{https://peopleoftheworld.net/graviton-pressure-theory/}.
\end{minipage}
\vfill
{\large Author: Shareef Ali Rashada\par}
{\normalsize shareef@peopleoftheworld.net\par}
\vspace{0.3cm}
{\normalsize Second edition, November 2026 \quad (first edition 12 June 2025)\par}
\end{titlepage}

\tableofcontents
\newpage

\section*{About This Paper}
\addcontentsline{toc}{section}{About This Paper}

This paper examines materialism as a theory of knowledge: the belief that physical matter is the only fundamental reality, and that whatever cannot be reduced to it is either derivative or illusory. It does so on materialism's own terms, without appeal to religion or the supernatural, by asking of each of its central explanations the question this framework asks of every explanation: does it name a mechanism, what touches what and what is exchanged, or does it give a name in place of one?\fn{rashadaPart1}

The paper stands on its own and needs no knowledge of Graviton Pressure Theory. Where it touches the framework's physics, it points to the part where that physics is developed. Its claims about purpose, will and meaning are philosophical claims, and are marked as such; the framework's physics papers make no claims about intent, and the full philosophical case is made in Part~2.\fn{rashadaPart2} This paper is where the two registers meet.

\statuskey

\section*{Introduction}
\addcontentsline{toc}{section}{Introduction}

Materialism has shaped modern science for centuries. It offered a clean, consistent and rational framework in which every phenomenon could be explained by particles, forces and blind interactions.

That framework is no longer held aloft by its own logic. The pressure against it is not mysticism or spiritual rebellion. It is \textbf{reason itself, pressing against materialism from within}. Where materialism treats the method of science as a finding about reality, it is a philosophy presented as physics, a dogma cloaked in empiricism.

This paper does not critique materialism from a religious standpoint. It examines the paradigm on \textit{its own terms}: its internal contradictions, its circular reasoning, and its growing reliance on unobservable constructs to patch its models.

We begin by separating the method of science from the metaphysics that has been attached to it. We then examine ``emergence'', and show where it explains and where it only names. We ask what moves a living body. We look at the ghost forces of cosmology, dark matter and dark energy, as they actually stand. We confront the fragmentation of knowledge, and the use of quantum measurement as an argument in either direction. And we propose a standard for knowledge that is higher than confirmation alone: \textbf{coherence, mechanism and fit with observation, together}.

This paper is a reckoning. It is not a retreat from science, but a maturation of it.

\section{Method and Metaphysics}
\label{sec:method}

A scientist at work studies matter: what it is made of, how it moves, what it does when it meets other matter. That is a method, and it has been extraordinarily successful. It is often called methodological naturalism: in the laboratory, look for natural causes, and test them.

Materialism is a different claim. It says that matter is \emph{all there is}: that mind, meaning and purpose are either arrangements of matter or illusions. That is not a method. It is a metaphysics, a claim about the whole of reality, and no experiment has ever tested it.

The two are routinely run together. ``We study matter'' becomes ``only matter exists''; ``we have not found a cause beyond matter'' becomes ``there is none''. The slide is rarely argued. It is assumed, and it is taught as though it were a result.

\argm{This paper's target is the slide, not the method.} The method of science stands. What does not stand is the habit of treating a successful method as a finding about everything, and then using that ``finding'' to rule out, in advance, whatever the method was not designed to detect.

\section{The Unpaid Debt}
\label{sec:debt}

Materialism promised simplicity: that all of reality could be reduced to matter, motion and mathematically defined interactions. It offered a seductive certainty: that anything not measurable was either derivative or illusory.

For a time, this worked. Materialism cut through superstition, challenged unexamined faiths and made physics the new authority of the rational. But as science has matured, the cost of the framework has shown itself.

Materialism has not been overthrown. Among professional philosophers it is still the majority view of the mind, but only just: in the 2020 PhilPapers Survey, 51.9 per cent accepted or leaned toward physicalism about the mind, and 32.1 per cent toward non-physicalism.\fn{bourget2023} That is not the profile of a settled result. It is the profile of a contested foundation.

What materialism has accumulated is a debt: explanations promised and not yet delivered. Consciousness is to be explained by emergence, someday. The motion of galaxies is to be explained by matter no one has detected. The acceleration of the universe is to be explained by an energy no one can describe. Each anomaly receives a placeholder, and each placeholder is a promissory note.

The cracks are not spiritual; they are \textbf{logical}. The symptoms are not theological; they are \textbf{epistemic}. A framework that once demanded evidence now asks for patience.

This paper does not call for a return to pre-scientific thinking. It calls for \textit{post-materialism}: a science that holds itself to the standard it once promised, of explanation, integration and causal coherence.

\section{Emergence: When It Explains and When It Names}
\label{sec:emergence}

Emergence is a legitimate idea, and science uses it well. Temperature emerges from the motion of molecules: no single molecule has a temperature, yet a gas does. Here emergence explains, because the mechanism is given. The molecules strike the wall of a thermometer, momentum and energy pass across the wall, and the reading measures the average energy of the strikes. What touches what, and what is exchanged, is named at every step.

Materialism insists that consciousness, a phenomenon defined by subjectivity, awareness and interiority, emerges from non-conscious matter in the same way. But here no mechanism is given. Neuroscience has traced, in great detail, which activity in the brain accompanies which experience. It has not said how any arrangement of matter becomes an experience: what, in the firing of a neuron, is the redness of red or the pain of a burn.

\argm{Used without a mechanism, ``emergence'' is not an explanation. It is a semantic shield}: a word placed where a mechanism should be, so that the question can be treated as answered. The reasoning beneath it runs: \textit{we do not know how consciousness comes from matter, but it must, because matter is all there is.} That is \textbf{circular}. It assumes the conclusion it is meant to support.
\begin{itemize}
  \item Consciousness is not a known property of matter, yet it is said to arise from particular arrangements of matter.
  \item Complexity alone is assumed to produce interiority, and no mechanism has been offered for how it would.
  \item The word is deployed not to advance understanding but to protect materialism from its greatest challenge: \textbf{the existence of mind as a primary reality}.
\end{itemize}

The philosopher David Chalmers made the problem precise. Even when every function in the vicinity of experience has been explained (perception, attention, report, the control of behaviour), a question remains: ``Why is the performance of these functions accompanied by experience?''\fn{chalmers1995} Why does all this processing not go on in the dark, with no inner feel at all? Chalmers' own answer was to propose that experience be taken as fundamental, as physics takes mass and charge as fundamental: not explained by anything simpler, but related to the rest of nature by laws.

A coherent epistemology must therefore either:
\begin{enumerate}
  \item provide a causal bridge from matter to mind, naming what touches what; or
  \item admit that consciousness may be as fundamental as mass or charge.
\end{enumerate}
Materialism, to preserve itself, chooses neither. It rebrands mystery as methodology, and hopes no one notices.

\section{What Moves a Living Body}
\label{sec:living}

The question of mind is often treated as a question about experience alone: what it is like to see, to feel, to think. There is a second question, and it is harder to set aside. Living things move themselves.

A stone falls when nothing holds it up. A person can decide to stand, and stand. A bird can turn against the wind. Part~7 of this framework draws the line exactly: the medium explains gravitational movement, but living bodies generate their own motion, and the laws that govern a falling stone do not decide what a living body will do next.\fn{rashadaPart7}

Physiology traces the chain from the brain to the muscle in detail: the signal in the motor cortex, the nerve, the release of calcium, the sliding of fibres that shortens the muscle. Every link after the first has a mechanism. The first link is the decision. Materialism answers that the decision \emph{is} the activity in the brain. But that is an identity asserted, not a mechanism shown. \argm{Ask what touches what when a choice moves an arm, and the materialist account names the arm, the nerve and the neuron, and gives a word, ``is'', where the choice meets the neuron.}

This is not an argument from ignorance. It is the interface test, applied as strictly to the mind as to gravity. If materialism is true, it owes an account of how a decision moves matter. Until it gives one, the moving living body stands as daily evidence that something other than matter in motion takes part in the world's causes.

\section{The Ghost Forces: Dark Matter and Dark Energy}
\label{sec:ghost}

Materialism often claims superiority on the grounds of predictive power. Yet the standard model of cosmology rests on two constructs that, together, make up about 95 per cent of the universe's mass and energy in that model, and neither has been detected directly.\fn{peebles2003}\fn{planck2020}

\paragraph{Dark matter.} Dark matter was proposed within Newtonian dynamics, not to rescue General Relativity. In 1933 Fritz Zwicky found that the galaxies of the Coma cluster move too fast to be held together by the matter that can be seen;\fn{zwicky1933} in the 1970s Vera Rubin and colleagues found that the outer parts of spiral galaxies orbit as fast as the inner parts, when the visible matter says they should orbit more slowly.\fn{rubin1980} Since then, several independent lines of evidence have been read the same way: the bending of light by clusters of galaxies, most famously in the Bullet Cluster;\fn{clowe2006} the pattern of fluctuations in the cosmic microwave background;\fn{planck2020} and the growth of structure in the universe.

The evidence that \emph{something} is missing is strong. What is missing is the thing itself. Decades of searches for a dark-matter particle, in detectors deep underground, have found nothing, and the most sensitive of them continue to set tighter limits.\fn{lz2023} And with each null result, the hypothesis moves: to lighter particles, weaker interactions, new candidates. A hypothesis that survives every failure by changing its parameters is falsifiable in principle and adjustable in practice.

The astronomer Stacy McGaugh, comparing the dark-matter paradigm with its main rival, wrote of the fixes added to models of galaxy formation to make them match what is seen: ``My fear is that feedback has become a modern version of the epicycle.''\fn{mcgaugh2015} The comparison is apt. In the Ptolemaic system, epicycles were not errors; they were fixes. They extended a flawed cosmology by patching over inconsistency, not by correcting its foundation.

\paragraph{Dark energy.} In 1998, observations of distant supernovae showed that the expansion of the universe is accelerating.\fn{riess1998}\fn{perlmutter1999} The cause was named dark energy. It has no mechanism: the simplest version, a constant energy of empty space, is smaller than the value quantum theory suggests by a factor so large that it is often called the worst prediction in physics. Recent surveys hint that dark energy may be changing over time, which, if confirmed, would mean the placeholder itself is not stable.\fn{desi2025}

\paragraph{What the ghost forces show.} Dark matter and dark energy are not frauds, and the observations behind them are real. They are names for what the observations require and the theory cannot supply. They are the dark sector of Part~3: explanations that account for an effect without naming what touches what.\fn{rashadaPart3} \hypm{The framework's own proposal, taken up in Part~9, is that the medium itself behaves differently far from a mass, where its flow is weakest; that proposal must meet the same data, and name its mechanism.} Whatever the answer proves to be, a science that holds itself to mechanism cannot count a placeholder as a discovery.

\section{The Neglect of Coherence}
\label{sec:coherence}

Materialism thrives on reduction. It treats complexity as a challenge to be taken apart, and the method has produced extraordinary results. But it has also produced a view of reality that is mechanically robust and philosophically hollow.

There are bridges between the sciences, in biophysics, neuroscience and the study of complex systems. What there is not is a shared account of what the world is, in which the findings of each fit together. Physics has no place for consciousness. Neuroscience studies the brain's activity and sets meaning aside. Cosmology describes the history of the universe without a mechanism for most of what it says the universe contains. Each field is expert within its walls. The result is not unified knowledge, but \textit{compartmentalised insight}: an epistemology that permits factual accuracy without existential intelligibility.

\begin{quote}
What does it mean that materialism has made reality more comprehensible, but less intelligible?
\end{quote}

It means that unity has been sacrificed for utility. It means that explanation has been severed from understanding. True understanding requires more than facts. It requires \textit{relationship among facts}. And until materialism treats the failure to integrate mind, matter, meaning and purpose as a failure, rather than as neutrality, it will continue to illuminate fragments while leaving the whole in darkness.

\section{Observation and the Measurement Problem}
\label{sec:measurement}

Materialism upholds objectivity as its central virtue: the removal of human subjectivity by attention to what is measurable and repeatable. That ideal is right in practice and incomplete in principle. All observation is made by a mind. All measurement begins in perception, and ends in someone's understanding of what was measured. There is no data without a frame, and no frame without a mind to hold it.

\begin{quote}
``Matter as primary'' is not an observation. It is an assumption.
\end{quote}

Quantum mechanics is often cited at this point, in both directions, and it needs to be cited accurately. John Wheeler wrote that ``no elementary phenomenon is a phenomenon until it is a registered (observed) phenomenon''.\fn{wheeler1983} But Wheeler meant registration by an irreversible act of amplification, such as the blackening of a grain of photographic emulsion or the click of a detector, and he said plainly that consciousness has nothing to do with the quantum process itself. Whether a mind plays any part in measurement is a question of interpretation, on which physicists disagree, and the data do not decide it. The claim that intention or expectation changes quantum outcomes is not supported.

What quantum mechanics does show is more useful to this paper's argument. Physics has a measurement problem, and has had one for a century. The theory describes a system evolving smoothly through many possibilities, and then, when it is measured, a single outcome. What happens in between is called ``collapse''. \argm{``Collapse'' is a name, not a mechanism: it does not say what touches what, or what is exchanged, when a possibility becomes a fact.} The rival interpretations (Copenhagen, many worlds, pilot waves, spontaneous collapse) agree on every prediction and disagree on what is happening. The foundation of modern physics, in other words, describes outcomes with great precision and leaves its central event unexplained.

That is the lesson. Not that mind controls matter, but that the most successful theory in physics has a gap at its centre where a mechanism should be, and that a science confident it has reduced everything to matter has not yet said what its own measurements are.

\section{Toward a Coherence-Based Epistemology}
\label{sec:epistemology}

We propose a higher standard for knowledge than confirmation alone. A theory must meet three tests together.
\begin{enumerate}
    \item \textbf{Coherence.} It must fit with what we know across domains, not only within its own: physics with life, life with mind, mind with the physics it uses to study the world.
    \item \textbf{Mechanism.} It must say what touches what, and what is exchanged. A name is not a mechanism, however precise the numbers it accounts for.
    \item \textbf{Fit with observation.} It must match what is measured, and say in advance what would show it wrong.
\end{enumerate}

Coherence alone is not enough. Philosophers have long noted that two systems of belief can each be perfectly coherent and still contradict each other; coherence must be anchored in observation. Observation alone is not enough either: a theory can match every measurement with placeholders, as the ghost forces show. Mechanism alone is not enough, since a mechanism can be invented. The three together are the standard, and they are the standard the rest of this framework applies to gravity.

\argm{Measured by that standard, materialism falls short on mind, on the living body, on the dark sector and on measurement, not because it has been refuted, but because at each it supplies a name where a mechanism is owed.}

A model of reality that meets all three must also make room for consciousness and causality, for purpose, will and structure, and must bring physics and mind into one account. Those are philosophical claims, made here as philosophy, and argued in full in Part~2. This is not a call for dogma. It is a call for \textbf{synthesis}: knowledge that does not fracture under magnification, nor collapse when viewed from afar.

\section{An Invitation to the Honest Scientist}
\label{sec:invitation}

What would it take for you to question your paradigm?
\begin{itemize}
  \item If an argument you could not fault led to consciousness as primary, would you follow it?
  \item If a framework explained more, with fewer placeholders and named mechanisms, would you examine it?
  \item If a spiritual framework offered greater explanatory coherence, would you be curious, or afraid?
\end{itemize}

This is not about religion. It is about \textbf{intellectual integrity}.

To question one's worldview is not betrayal; it is responsibility. Science at its highest is not the defence of consensus. It is the pursuit of coherence through evidence, reason and lived alignment.

We are not asking you to believe.\\
We are asking you to \textit{think}: without reflex, without fear, and without allegiance to paradigms that have outlived their explanatory power.

\section*{Conclusion: The Era of Assumption Is Ending}
\addcontentsline{toc}{section}{Conclusion: The Era of Assumption Is Ending}

Materialism is not being undone by theology, mysticism or resistance to progress. It is being undone by its own debts: the explanations it has promised and not delivered.

Its watchwords, \textit{emergence, randomness, reduction}, have not explained the world we actually live in: a world of structure, consciousness, intention and coherence. At each point where the explanation matters most, they give a name where a mechanism is owed.

We are not replacing science with mysticism. We are holding science to its own standard.

A new epistemology is rising, one that sees:
\begin{itemize}
  \item mind not as an accident, but as an anchor;
  \item meaning not as illusion, but as structural;
  \item structure not as a byproduct of randomness, but as the very shape of existence.
\end{itemize}

Its time has come.

\section*{Changes in the Second Edition}
\addcontentsline{toc}{section}{Changes in the Second Edition}

This list records what changed from the first edition (12 June 2025) and why. Changes to the argument alter what the paper claims or how it supports a claim; clarifications improve wording, sources or presentation and leave the argument as it was. The paper's title, voice, length and central case are kept.

\subsection*{Changes to the argument}
\begin{itemize}
  \item New Section~12.1 separates the method of science (methodological naturalism) from materialism as a metaphysics, and makes the slide from one to the other the paper's target. The first edition treated the two as one, which left the argument open to the reply that it attacks a strawman.
  \item ``Collapse'' is scoped. Materialism is shown as a contested foundation carrying an unpaid debt of promised explanations, with the 2020 PhilPapers Survey (51.9 per cent physicalism, 32.1 per cent non-physicalism about the mind). The title is kept.
  \item The critique of emergence now applies the interface test: emergence explains where a mechanism is given (temperature from molecular motion) and only names where none is given (consciousness).
  \item New Section~12.4, on what moves a living body: living things generate their own motion (Part~7), and materialism owes an account of how a decision moves matter.
  \item The section on dark matter and dark energy is corrected and strengthened. Dark matter was proposed within Newtonian dynamics (Zwicky 1933; Rubin et al.\ 1980), not to preserve General Relativity, and several independent lines of evidence point to something missing. The critique is restated as the absence of detection and mechanism, and the adjustability of the hypothesis after each null search. Dark energy's lack of mechanism and recent hints that it changes over time are added.
  \item The section on objectivity is recast. The first edition's claim that ``intention and expectation influence quantum states'' is withdrawn as unsupported, and Wheeler is cited for what he wrote, including his statement that consciousness has nothing to do with the quantum process. The section now makes the stronger point: the measurement problem is unsolved, and ``collapse'' is a name, not a mechanism.
  \item The coherence standard becomes three tests together (coherence, mechanism and fit with observation), answering the objection that coherence alone cannot choose between rival coherent systems.
  \item The question ``If AI systems consistently rejected materialism, would you listen?'' is replaced by questions about following an argument and examining a framework with fewer placeholders.
\end{itemize}

\subsection*{Clarifications}
\begin{itemize}
  \item The quotation attributed to Chalmers (``consciousness is not explained by functional organization---it is something extra'') could not be found in the cited paper and is replaced by his own question, with his proposal to take experience as fundamental given in paraphrase.
  \item The quotation attributed to McGaugh (``We invented dark matter \ldots\ It's the 21st century's epicycle'') could not be found in the cited paper and is replaced by his actual sentence about feedback as ``a modern version of the epicycle''.
  \item Claims about purpose, will and meaning are identified as philosophical claims, with the full case in Part~2.
  \item New sources: Bourget and Chalmers (2023), Zwicky (1933), Rubin et al.\ (1980), Clowe et al.\ (2006), Planck (2020), the LZ collaboration (2023), Riess et al.\ (1998), Perlmutter et al.\ (1999) and DESI (2025). Wheeler (1983) is cited as his essay ``Law without Law''.
  \item An ``About This Paper'' section and this list are added, and status marks are shown on the framework's own claims (see the key at the start).
\end{itemize}

\newpage
\bibliographystyle{plainnat}
\bibliography{part12}

\end{document}
